% GENERATED FILE. Edit canonical lessons, then run npm run generate:books.
% canonical-source-hash: fnv1a64:36e5f8765cb9e6f9
% canonical-lessons: MR-C425-taraju, MR-C425-mapne, MR-C425-assal, MR-C425-banavat, MR-C425-nakkal

\chapter{Weights and Measures}
\label{ch:mr-weights-and-measures}
\hlchaptermodality{425}

\begin{chapteropening}
\textbf{By the end of this chapter:} \emph{I can talk about scales and measuring, and say what is genuine, fake or a copy.}

Complete the last lesson of chapter 425: I can talk about scales and measuring, and say what is genuine, fake or a copy.
\end{chapteropening}

\section[\texorpdfstring{\mr{तराजू} (tarājū)}{तराजू (tarājū)}]{\mr{तराजू} (tarājū) — a pair of scales}

\label{lesson:MR-C425-taraju}



\begin{quote}
Before the new one: say the Marathi for a broker, then the Marathi for a bill.
\end{quote}

\subsection*{You'll want to know: \mr{तराजू}}
\textbf{\mr{तराजू}} — \emph{tarājū} — \textquotedblleft{}a pair of scales\textquotedblright{}.

\begin{quote}
\mr{भाजीवाला} \mr{तराजूने} \mr{कांद्यांचं} \mr{वजन} \mr{करतो}.
\end{quote}

\begin{quote}
\emph{bhājīvālā tarājūne kāndyāñcaṁ vajan karto.}
\end{quote}

\begin{quote}
The vegetable seller weighs the onions on the scales.
\end{quote}

Masculine.

A masculine noun; the plural is the same.

\begin{cousinweb}
\textbf{\mr{तराजू}} (\emph{tarājū}) is a loan usually traced to Persian \emph{tarāzū}, \textquotedblleft{}a balance, a pair of scales\textquotedblright{}.
\end{cousinweb}

\begin{grammarlens}[title={What you've built}]
One more word for getting around.
\end{grammarlens}

\subsection*{Your turn}
\begin{itemize}
\raggedright
  \item \emph{Say it:} \emph{tarājū}
  \item \emph{Say it:} \emph{tarājū}, once more
  \item \emph{Say it:} \emph{bhājīvālā tarājūne kāndyāñcaṁ vajan karto.} — The vegetable seller weighs the onions on the scales.
  \item \emph{Recall:} say the Marathi for a broker, then the Marathi for a bill, then say \emph{tarājū} again
\end{itemize}

\begin{tcolorbox}[breakable,colback=teal!4,colframe=teal!35!black,title={Before you move on}]
What does \mr{तराजू} mean? (\textquotedblleft{}A pair of scales\textquotedblright{}.) Say it once more.
\end{tcolorbox}
\section[\texorpdfstring{\mr{मापणे} (māpṇe)}{मापणे (māpṇe)}]{\mr{मापणे} (māpṇe) — to measure}

\label{lesson:MR-C425-mapne}



\begin{quote}
Before the new one: say the Marathi for a bill, then the Marathi for a pair of scales.
\end{quote}

\subsection*{You'll want to know: \mr{मापणे}}
\textbf{\mr{मापणे}} — \emph{māpṇe} — \textquotedblleft{}to measure\textquotedblright{}.

\begin{quote}
\mr{मी} \mr{खोलीची} \mr{लांबी} \mr{मापली}.
\end{quote}

\begin{quote}
\emph{mī kholīcī lāmbī māplī.}
\end{quote}

\begin{quote}
I measured the length of the room.
\end{quote}

Past, with -\mr{ने} (he …): \mr{त्याने} \mr{मापलं}, \emph{tyāne māplaṁ}.

The thing measured is the object: \mr{लांबी} \mr{मापली}. The noun is \mr{माप}: \mr{माप} \mr{घेणे}, to take measurements.

\begin{cousinweb}
\textbf{\mr{मापणे}} (\emph{māpṇe}) is inherited: CDIAL lists Marathi \emph{māpṇẽ}, \textquotedblleft{}to measure\textquotedblright{}, under \emph{*māpyate}, \textquotedblleft{}is measured\textquotedblright{} (CDIAL 10054).
\end{cousinweb}

\begin{grammarlens}[title={What you've built}]
One more word for getting around.
\end{grammarlens}

\subsection*{Your turn}
\begin{itemize}
\raggedright
  \item \emph{Say it:} \emph{māpṇe}
  \item \emph{Say it:} \emph{māpṇe}, once more
  \item \emph{Say it:} \emph{mī kholīcī lāmbī māplī.} — I measured the length of the room.
  \item \emph{Recall:} say the Marathi for a bill, then the Marathi for a pair of scales, then say \emph{māpṇe} again
\end{itemize}

\begin{tcolorbox}[breakable,colback=teal!4,colframe=teal!35!black,title={Before you move on}]
What does \mr{मापणे} mean? (\textquotedblleft{}To measure\textquotedblright{}.) Say it once more.
\end{tcolorbox}
\section[\texorpdfstring{\mr{अस्सल} (assal)}{अस्सल (assal)}]{\mr{अस्सल} (assal) — genuine}

\label{lesson:MR-C425-assal}



\begin{quote}
Before the new one: say the Marathi for a pair of scales, then the Marathi for to measure.
\end{quote}

\subsection*{You'll want to know: \mr{अस्सल}}
\textbf{\mr{अस्सल}} — \emph{assal} — \textquotedblleft{}genuine\textquotedblright{}.

\begin{quote}
\mr{हे} \mr{अस्सल} \mr{सोनं} \mr{आहे}.
\end{quote}

\begin{quote}
\emph{he assal sonaṁ āhe.}
\end{quote}

\begin{quote}
This is genuine gold.
\end{quote}

It does not change: \mr{अस्सल} \mr{रेशीम}, \mr{अस्सल} \mr{पुणेरी}.

\begin{cousinweb}
\textbf{\mr{अस्सल}} (\emph{assal}) is a loan through Persian, usually traced to Arabic \emph{aṣl}, \textquotedblleft{}root, origin\textquotedblright{}: something true to its origin.
\end{cousinweb}

\begin{grammarlens}[title={What you've built}]
One more word for getting around.
\end{grammarlens}

\subsection*{Your turn}
\begin{itemize}
\raggedright
  \item \emph{Say it:} \emph{assal}
  \item \emph{Say it:} \emph{assal}, once more
  \item \emph{Say it:} \emph{he assal sonaṁ āhe.} — This is genuine gold.
  \item \emph{Recall:} say the Marathi for a pair of scales, then the Marathi for to measure, then say \emph{assal} again
\end{itemize}

\begin{tcolorbox}[breakable,colback=teal!4,colframe=teal!35!black,title={Before you move on}]
What does \mr{अस्सल} mean? (\textquotedblleft{}Genuine\textquotedblright{}.) Say it once more.
\end{tcolorbox}
\section[\texorpdfstring{\mr{बनावट} (banāvaṭ)}{बनावट (banāvaṭ)}]{\mr{बनावट} (banāvaṭ) — fake; counterfeit}

\label{lesson:MR-C425-banavat}



\begin{quote}
Before the new one: say the Marathi for to measure, then the Marathi for genuine.
\end{quote}

\subsection*{You'll want to know: \mr{बनावट}}
\textbf{\mr{बनावट}} — \emph{banāvaṭ} — \textquotedblleft{}fake; counterfeit\textquotedblright{}.

\begin{quote}
\mr{पोलिसांनी} \mr{बनावट} \mr{नोटा} \mr{पकडल्या}.
\end{quote}

\begin{quote}
\emph{polisānnī banāvaṭ noṭā pakḍalyā.}
\end{quote}

\begin{quote}
The police seized fake banknotes.
\end{quote}

It does not change: \mr{बनावट} \mr{नोट}, \mr{बनावट} \mr{सही}.

\begin{cousinweb}
\textbf{\mr{बनावट}} (\emph{banāvaṭ}) is a loan, usually traced to Hindustani \emph{banāvaṭ}, \textquotedblleft{}make, something put together\textquotedblright{}.
\end{cousinweb}

\begin{grammarlens}[title={What you've built}]
One more word for getting around.
\end{grammarlens}

\subsection*{Your turn}
\begin{itemize}
\raggedright
  \item \emph{Say it:} \emph{banāvaṭ}
  \item \emph{Say it:} \emph{banāvaṭ}, once more
  \item \emph{Say it:} \emph{polisānnī banāvaṭ noṭā pakḍalyā.} — The police seized fake banknotes.
  \item \emph{Recall:} say the Marathi for to measure, then the Marathi for genuine, then say \emph{banāvaṭ} again
\end{itemize}

\begin{tcolorbox}[breakable,colback=teal!4,colframe=teal!35!black,title={Before you move on}]
What does \mr{बनावट} mean? (\textquotedblleft{}Fake; counterfeit\textquotedblright{}.) Say it once more.
\end{tcolorbox}
\section[\texorpdfstring{\mr{नक्कल} (nakkal)}{नक्कल (nakkal)}]{\mr{नक्कल} (nakkal) — a copy; an imitation}

\label{lesson:MR-C425-nakkal}



\begin{quote}
Before the new one: say the Marathi for genuine, then the Marathi for fake.
\end{quote}

\subsection*{You'll want to know: \mr{नक्कल}}
\textbf{\mr{नक्कल}} — \emph{nakkal} — \textquotedblleft{}a copy; an imitation\textquotedblright{}.

\begin{quote}
\mr{हा} \mr{फोटो} \mr{मूळ} \mr{चित्राची} \mr{नक्कल} \mr{आहे}.
\end{quote}

\begin{quote}
\emph{hā phoṭo mūḷ citrācī nakkal āhe.}
\end{quote}

\begin{quote}
This photo is a copy of the original painting.
\end{quote}

Feminine.

\mr{नक्कल} \mr{करणे}, to copy or to mimic someone.

\begin{cousinweb}
\textbf{\mr{नक्कल}} (\emph{nakkal}) is a loan through Persian, usually traced to Arabic \emph{naql}, \textquotedblleft{}a transfer, a copy\textquotedblright{}.
\end{cousinweb}

\begin{grammarlens}[title={What you've built}]
One more word for getting around.
\end{grammarlens}

\subsection*{Your turn}
\begin{itemize}
\raggedright
  \item \emph{Say it:} \emph{nakkal}
  \item \emph{Say it:} \emph{nakkal}, once more
  \item \emph{Say it:} \emph{hā phoṭo mūḷ citrācī nakkal āhe.} — This photo is a copy of the original painting.
  \item \emph{Recall:} say the Marathi for genuine, then the Marathi for fake, then say \emph{nakkal} again
\end{itemize}

\begin{tcolorbox}[breakable,colback=teal!4,colframe=teal!35!black,title={Before you move on}]
What does \mr{नक्कल} mean? (\textquotedblleft{}A copy; an imitation\textquotedblright{}.) Say it once more.
\end{tcolorbox}
